% META: {"id":"1SPE-PRODSCAL-CO-047","chapitre":"1SPE-PRODUIT-SCALAIRE","type_objet":"corrige","exercice_id":"1SPE-PRODSCAL-EX-047","status":"generated"}
% BEGIN-VERIFY
% from sympy import Matrix, symbols, cos, sin, trigsimp, expand, pi
% b, c = symbols('b c', positive=True)
% theta = symbols('theta', real=True)
% # Construction du troisième côté dans un repère orthonormé.
% AB = Matrix([c, 0])
% AC = Matrix([b*cos(theta), b*sin(theta)])
% BC = AC - AB
% a_sq = BC.dot(BC)
% assert AB.dot(AB) == c**2
% assert trigsimp(AC.dot(AC) - b**2) == 0
% assert AB.dot(AC) == b*c*cos(theta)
% assert trigsimp(a_sq - (b**2+c**2-2*b*c*cos(theta))) == 0
% # Cas particuliers complémentaires ; ils ne remplacent pas l'identité générale.
% assert trigsimp(a_sq.subs(theta, pi/2) - (b**2+c**2)) == 0
% assert expand(a_sq.subs(theta, pi/3) - (b**2+c**2-b*c)) == 0
% END-VERIFY
\begin{corrige}{1SPE-PRODSCAL-EX-047}
$\overrightarrow{BC} = \overrightarrow{AC} - \overrightarrow{AB}$, donc :
\begin{align*}
BC^2 &= \|\overrightarrow{BC}\|^2 = \|\overrightarrow{AC} - \overrightarrow{AB}\|^2 \\
&= \|\overrightarrow{AC}\|^2 - 2\,\overrightarrow{AB} \cdot \overrightarrow{AC} + \|\overrightarrow{AB}\|^2 \\
&= AC^2 + AB^2 - 2\,AB \times AC \times \cos\widehat{BAC}.
\end{align*}
\end{corrige}
